% Ported verbatim (prose unchanged) from docs/SECTION_3_DRAFT.md.
\section{Attribution Methodology}
\label{sec:methodology}

The central question of this paper --- whether conversation-graph machinery retains measurable value once the utterance encoder is fine-tuned --- is an attribution question, and attribution questions are notoriously vulnerable to confounds that ordinary ablation studies do not control. When a new architectural component is added to a model and the ensemble is retrained from scratch, any difference in final performance conflates the component's representational contribution with incidental effects of its presence: additional randomly initialized parameters that must be optimized before the model recovers its baseline competence, altered optimization dynamics, and increased overfitting capacity on modestly sized corpora. In our own preliminary experiments (Section~\ref{sec:meld-mechanism}), these incidental effects were not small: the train--validation gap grew monotonically with component count, from $+0.129$ for the fusion baseline to $+0.225$ for the full model, and configurations with more components underperformed configurations with fewer even before any question of component quality arose. A methodology that cannot distinguish ``the component adds nothing'' from ``the component's presence disturbed training'' cannot answer our question. This section describes the design that makes the distinction, in four parts: the fine-tuned foundation on which all configurations rest (\ref{sec:meld-foundation}), the residual construction under which every component provably initializes as an identity (\ref{sec:residual-construction}), the components under test (\ref{sec:components-under-test}), and the experimental protocol (\ref{sec:protocol}).

\subsection{The Fine-Tuned Textual Foundation}
\label{sec:meld-foundation}

All configurations share a single textual foundation per corpus: a RoBERTa-base encoder adapted with Low-Rank Adaptation (LoRA)~\cite{hu2022lora} applied to the query and value projections of every self-attention layer (rank $r{=}8$, scaling $\alpha{=}16$, adapter dropout $0.05$), with the base weights frozen. The input for utterance $t$ is a context-window construction: the $k$ most recent preceding utterances of the same dialogue, each prefixed with its speaker's name, concatenated with separator tokens and followed by the current utterance, truncated oldest-first to $256$ tokens while always preserving the current utterance intact ($k{=}8$ unless stated otherwise; Section~\ref{sec:k-sweep} sweeps $k$). No future context is visible. The utterance representation is the attention-mask-weighted mean of the second-to-last layer's hidden states, a choice selected once on validation data and locked; on our features this pooling outperformed both the final layer's mean and the untuned CLS position, consistent with prior observations that RoBERTa's sentence-level CLS representation is weak without task-specific tuning.

Two implementation findings from constructing this foundation are worth recording because they contradict reasonable defaults. First, RoBERTa-base's \texttt{pooler\_output} head is randomly initialized at load time and is therefore process-dependent: features cached across process restarts occupy different random projection spaces, which in an early pipeline of ours produced training and test features in inconsistent spaces and a catastrophic, difficult-to-diagnose collapse. All cached representations in this work are verified process-independent (twenty utterances re-encoded in a fresh process must match the cache to within $10^{-3}$) before any training run consumes them. Second, training-time logit adjustment~\cite{menon2021logitadjustment} --- subtracting $\tau \log \pi_c$ from the logits inside the loss --- failed in our setting in a specific and diagnosable way: under the trained-adjustment variant, the two rarest classes were ranked last among seven classes by the raw logits for $82\%$ and $69\%$ of their test instances respectively, and no inference-time correction could recover them. Plain cross-entropy training with post-hoc logit adjustment at inference (the estimator's alternative form) moved those ranks to the top three-to-four and left both classes with nonzero F1 in every seed. We therefore train all classifiers with unadjusted cross-entropy and, where a validation-selected adjustment temperature qualifies under a fixed rule, apply adjustment only at inference; on IEMOCAP no temperature qualified and raw logits are used throughout. We report this negative result because trained logit adjustment is widely cited and its silent failure mode --- excellent aggregate metrics with structurally dead rare classes --- is easy to miss.

The foundation is trained once per corpus per seed and then frozen. On MELD it attains a test weighted F1 of $0.6403 \pm 0.0045$ across three seeds --- exceeding, for calibration, the published performance of several complete multimodal graph systems, including our own prior SocialArcNet (published: $0.62$ weighted F1)~\cite{socialarcnet2026}. This is the first substantive observation of the paper: the foundation that graph architectures are asked to improve upon has moved.

\subsection{The Residual ``Do-No-Harm'' Construction}
\label{sec:residual-construction}

Let $z^{\text{text}}_t$ denote the frozen foundation's logits for utterance $t$, cached alongside its pooled embedding. Every configuration in this study produces its prediction as
\begin{equation}
z_t \;=\; z^{\text{text}}_t \;+\; W_{\text{out}}\, g_t,
\label{eq:residual}
\end{equation}
where $g_t$ is the readout of the configuration's architectural stack --- multimodal fusion, speaker-state recurrence, relational memory, auxiliary heads, in whatever combination the configuration enables --- and $W_{\text{out}}$ is initialized to zero. Additionally, in the fusion projector that combines the textual embedding with the (frozen, cached) acoustic and visual features, the blocks acting on the acoustic and visual inputs are zero-initialized, so that non-textual modalities enter only as learned corrections. The consequence is a property we verify rather than assume: at initialization, the predictions of every configuration are exactly the predictions of the frozen foundation, for every input. This equality-at-init property is enforced by an automated test parametrized over all component-flag combinations and executed against real cached data before any training run of a campaign is permitted to launch.

The construction converts the attribution question into a well-posed one. Because every configuration begins at the foundation's performance and gradient descent is free to leave $W_{\text{out}}$ at zero, a configuration ends below the foundation only if training actively moved it there, and ends above it only if the stack found signal the foundation lacks. The incidental-presence confound is eliminated by design: there is no baseline competence to re-learn, because the baseline is the starting point. The measured quantity --- the paired, per-seed difference between a configuration and the foundation, or between two configurations differing in a single component --- is attributable to the component under a construction argument rather than a ceteris-paribus hope. The design has precedent in spirit (zero-initialized residual branches appear in normalization-free networks and in adapter literature), but we are not aware of its prior use as an attribution instrument for conversation-level ERC architecture, and Section~\ref{sec:meld} shows that it materially changes conclusions: the same components that appear harmful under scratch retraining ($-0.0129$ versus baseline) appear merely inert under residual attribution ($-0.0048$), and the difference between those two readings is precisely the confound the design removes.

Two honest limitations of the construction are recorded here and revisited in Section~\ref{sec:analysis}. First, zero-initialization is a soft prior toward the foundation; a component whose benefit requires wholesale reorganization of the shared pathway could in principle be disadvantaged. We mitigate this by also reporting the scratch-retrained results (Section~\ref{sec:meld-finetuned}), by verifying that gradients flow to all stack parameters at initialization, and by observing that the identical apparatus does detect positive contributions where they exist --- on frozen features (Section~\ref{sec:meld-frozen}) the graph stack's paired gain is $+0.0356 \pm 0.0106$ under the same construction. The instrument is demonstrably capable of measuring a positive effect; it simply measures none on fine-tuned foundations. Second, the foundation checkpoint is fixed per corpus (seed $42$) while downstream seeds vary head initialization and data order; foundation-seed variance is therefore absorbed into the anchor rather than the deltas, a choice that narrows the comparison's variance at the cost of one untested interaction, which we accept and disclose.

\subsection{Components Under Test}
\label{sec:components-under-test}

The stack instantiates, behind independent configuration flags, the principal mechanisms of the graph-ERC literature together with two proposals of our own. The speaker-state pathway maintains a hidden state per speaker (dimension $256$), updated at each utterance by a graph-attention message-passing step over all speakers' previous states followed by a gated recurrent update conditioned on the fused utterance embedding --- the standard recurrent-graph pattern of contextual ERC models, including our prior work. The relational-memory pathway, our first proposal, maintains a recurrent state per unordered speaker pair (dimension $128$, one shared update cell across pairs, zero-initialized per dialogue), updated whenever either member speaks and injected into message passing through edge-conditioned attention in the GATv2 form; the classification readout concatenates the speaker's state with the mean of its incident edge states, making relationships first-class in the output pathway. The shift-aware objective, our second proposal, adds a binary auxiliary head predicting whether the current utterance's emotion differs from the same speaker's previous emotion (first utterances masked), trained with class-weighted binary cross-entropy at fixed weight $\lambda{=}0.5$ against the primary loss; the premise --- that rare emotions arrive disproportionately as shifts --- is verified against corpus statistics before use (neutral shift rate $0.39$--$0.43$ across MELD's splits vs.\ fear/disgust/surprise at $1.5$--$2.1\times$ higher, no exception across train/dev/test; on IEMOCAP the same check is only partially confirmed --- Section~\ref{sec:dissociation} reports both in full). Temporal self-attention pooling replaces mean-pooling over frame-level visual and step-level acoustic features with a single attention block whose initialization provably reproduces mean-pooling, the same do-no-harm principle applied one level down. Acoustic and visual features come from frozen Wav2Vec~2.0 and MViTv2 backbones respectively; their linear-probe ceilings on MELD ($0.41$ and $0.33$ weighted F1) are reported so that the reader can calibrate how much signal the non-textual modalities offer any architecture.

\subsection{Protocol}
\label{sec:protocol}

All substantive analysis choices were fixed before the results they govern were known, and the full history is public. Hypotheses, gates, and decision rules were written into the repository in advance of the experiments that test them; in the one case where a hypothesis concerned a dataset we did not yet possess (the IEMOCAP relational-memory trial, Section~\ref{sec:iemocap}), the hypothesis and its reading rule were committed before the corpus was downloaded, and the commit hashes are reported alongside the results. Training recipes --- optimizer, schedule, losses, selection metric, early stopping, seeds $\{42, 1337, 2024\}$, extended to seven seeds where a null claim required tighter error bars --- are recorded in a locked recipe document; a post-hoc audit of every run's committed source confirmed no drift between corpora. Model selection and early stopping use validation macro F1 throughout (chosen after observing that weighted-F1 selection structurally abandons rare classes; the exception is documented at its origin), while weighted F1 remains the primary reported metric for comparability. With three to seven seeds we report means, standard deviations, and paired per-seed differences, and we do not perform hypothesis tests that would imply statistical power we do not have. Every result table in this paper is regenerated from stored per-run metrics by a single script; the repository, including the complete decision history, failed intermediate designs, and this methodology's own development, is available at \url{https://github.com/multi-modal-rtm/rapport}.
